\subsubsection{Convolution of Distributions}\label{subsec:distribution-convolution}

Let $f$ and $g$ be continuous functions on $\mathbb{R}^m$ which decay rapidly at infinity. We then have the following identity:
\begin{equation}\label{eq:convolution-test-pairing-integral-identity}
\begin{aligned}
(f*g,\phi)&=\int_{\mathbb{R}^m}(f*g)(\bm{x})\phi(\bm{x})\,\mathrm{d}\bm{x}\\
&=\int_{\mathbb{R}^m}\int_{\mathbb{R}^m}f(\bm{x}-\bm{y})g(\bm{y})\phi(\bm{x})\,\mathrm{d}\bm{x}\,\mathrm{d}\bm{y}\\
&=\int_{\mathbb{R}^m}\int_{\mathbb{R}^m}f(\bm{x})g(\bm{y})\phi(\bm{x}+\bm{y})\,\mathrm{d}\bm{x}\,\mathrm{d}\bm{y}.
\end{aligned}
\end{equation}
This identity is used as the definition of the convolution of two distributions.

\begingroup
\RenewDocumentEnvironment{Definition}{g}
{
    \refstepcounter{statement}
    {\sffamily\bfseries\color{PageColor}\fontsize{10.95}{13.2}\selectfont
     \textquotedbl Definition\textquotedbl~\thestatement
     \IfValueT{#1}{{\mdseries\ (#1)}}\textbf{:}}\quad
    \kaishu
}
{\par}
\begin{Definition}
\label{def:distribution-convolution-formal}
Let $f,g\in\mathcal{D}^{\prime}(\mathbb{R}^m)$. Then the convolution of $f$ and $g$ is defined by
\begin{equation}\label{eq:distribution-convolution-formal-definition}
(f*g,\phi)=(f(\bm{x})g(\bm{y}),\phi(\bm{x}+\bm{y})).
\end{equation}

\end{Definition}
\endgroup

The quotes are meant to draw attention to the fact that this does not make any sense. We defined the direct product $f(\bm{x})g(\bm{y})$ as an element of $\mathcal{D}^{\prime}(\mathbb{R}^{2m})$, but $\phi(\bm{x}+\bm{y})$ is not in $\mathcal{D}(\mathbb{R}^{2m})$; it does not have compact support. Indeed, the convolution of arbitrary distributions cannot be defined in a rational manner. There are, however, special cases where a meaning can be given to \eqref{eq:distribution-convolution-formal-definition}. In the simplest such case, the support of $\phi(\bm{x}+\bm{y})$ has a compact intersection with the support of $f(\bm{x})g(\bm{y})$. If this is the case, we may replace $\phi(\bm{x}+\bm{y})$ by any test function which agrees with $\phi(\bm{x}+\bm{y})$ in a neighborhood of $\operatorname{supp}(f(\bm{x})g(\bm{y}))$. In particular, \eqref{eq:distribution-convolution-formal-definition} is meaningful under either of the following conditions:
\begin{enumerate}
\item Either $f$ or $g$ has compact support.
\item In one dimension, the supports of $f$ and $g$ are bounded from the same side (e.g., $f=0$ for $x<a$ and $g=0$ for $x<b$).
\end{enumerate}

From the corresponding properties of the direct product, it follows that convolution is commutative and associative where it is defined.

Let us consider some special cases:

\begin{enumerate}\item We have
\begin{equation}\label{eq:delta-convolution-identity-pairing}
\begin{aligned}
(\delta*f,\phi)&=(\delta(\bm{x})f(\bm{y}),\phi(\bm{x}+\bm{y}))\\
&=(f(\bm{y}),(\delta(\bm{x}),\phi(\bm{x}+\bm{y})))\\
&=(f(\bm{y}),\phi(\bm{y}))\\
&=(f,\phi),
\end{aligned}
\end{equation}
i.e., $\delta*f=f$.
\item Let us consider $f*\psi$, where $\psi\in\mathcal{D}(\mathbb{R}^m)$. We have
\begin{equation}\label{eq:distribution-test-convolution-smooth-pairing}
\begin{aligned}
(f*\psi,\phi)&=(f(\bm{x})\psi(\bm{y}),\phi(\bm{x}+\bm{y}))\\
&=\left(f(\bm{x}),\int_{\mathbb{R}^m}\psi(\bm{y})\phi(\bm{x}+\bm{y})\,\mathrm{d}\bm{y}\right)\\
&=\left(f(\bm{x}),\int_{\mathbb{R}^m}\psi(\bm{y}-\bm{x})\phi(\bm{y})\,\mathrm{d}\bm{y}\right)\\
&=\int_{\mathbb{R}^m}(f(\bm{x}),\psi(\bm{y}-\bm{x}))\phi(\bm{y})\,\mathrm{d}\bm{y}.
\end{aligned}
\end{equation}
In the last step, we have used the continuity of the functional $f$ to take it under the integral; see Problem~\ref{ex:test-function-integral-riemann-sums}. Hence $f*\psi(\bm{y})$ is equal to the function $(f(\bm{x}),\psi(\bm{y}-\bm{x}))$. This function is of class $C^\infty$, and if $f$ has compact support, it is a test function.
\end{enumerate}

We next consider differentiation of a convolution. By definition, we have
\begin{equation}\label{eq:distribution-convolution-derivative-pairing}
\begin{aligned}
(D^\alpha(f*g),\phi)&=(-1)^{|\alpha|}(f*g,D^\alpha\phi)\\
&=(-1)^{|\alpha|}(g(\bm{y}),(f(\bm{x}),D^\alpha\phi(\bm{x}+\bm{y})))\\
&=(g(\bm{y}),(D^\alpha f(\bm{x}),\phi(\bm{x}+\bm{y})))\\
&=(D^\alpha f*g,\phi).
\end{aligned}
\end{equation}
Thus $D^\alpha(f*g)=D^\alpha f*g$, and using commutativity, we also find $D^\alpha(f*g)=f*D^\alpha g$. A convolution is differentiated by differentiating either one of the factors.

The following lemma expresses continuity of the operation of convolution.

\begin{Lemma}
\label{lem:distribution-convolution-continuity}
Assume that $f_n\to f$ in $\mathcal{D}^{\prime}(\mathbb{R}^m)$ and that one of the following holds:
\begin{enumerate}
\item The supports of all the $f_n$ are contained in a common compact set;
\item $g$ has compact support;
\item $m=1$ and the supports of the $f_n$ and of $g$ are bounded on the same side, independently of $n$.
\end{enumerate}
Then $f_n*g\to f*g$ in $\mathcal{D}^{\prime}(\mathbb{R}^m)$.
\end{Lemma}

The proof is left as an exercise (cf. Problem~\ref{ex:distribution-convolution-continuity}). A consequence is the following theorem.

\begin{Theorem}
\label{thm:test-functions-dense-distributions}
$\mathcal{D}(\mathbb{R}^m)$ is dense in $\mathcal{D}^{\prime}(\mathbb{R}^m)$.
\end{Theorem}

\begin{Proof}
We first show that distributions of compact support are dense. To see this, simply let $e_n$ be a test function which equals 1 on the set $\{|\bm{x}|\leq n\}$. Then $e_nf\to f$ for every $f\in\mathcal{D}^{\prime}(\mathbb{R}^m)$, and the support of $e_nf$ is contained in the support of $e_n$, hence compact.

It therefore suffices to show that distributions of compact support are limits of test functions. Let $f$ be a distribution of compact support, and let $\phi_n$ be a delta-convergent sequence of test functions; we may for example choose the sequence $\phi_n=C(1/n)^{-1}\phi_{0,1/n}$, where $\phi_{0,\epsilon}$ and $C(\epsilon)$ are defined by \eqref{eq:smooth-compact-bump-function} and \eqref{eq:test-function-bump-normalization}. Then $\phi_n*f$ is a test function, and by the previous lemma $\phi_n*f$ converges to $\delta*f=f$.
\end{Proof}
